% !TEX root = canbeincluded.tex
% LTeX: language = nl
% !TeX spellcheck = nl_NL
\newcommand{\myClass}{article}
\newcommand{\cursuscode}{\Cpp}
\newcommand{\titel}{Objectgeoriënteerd Programmeren in \Cpp{}}
\newcommand{\documenttype}{Notities}
\newcommand{\huidigeversie}{1.2}
\newcommand{\auteurs}{
	J.Z.M. Broeders
}
\include{style_nl_NL}

\usepackage{needspace}
\usepackage{hwemoji}
\usepackage{xfrac}
% \usepackage{comment}
\newcommand{\myfrac}[2]{\raisebox{0.35ex}{\sfrac{#1}{#2}}}

\newcommand{\progs}{progs}
\newcommand{\progsurl}{https://bitbucket.org/HR_ELEKTRO/cppprog/src/master/progs}
\include{styleCode}

% Commando \cpplstinputfragment leest code uit (een deel van) een bestand en plaatst dit in een codefragment.
% #1 = (optional) extra key=value paren voor lstinputlisting b.v.
%      style = linenumbers
%      linerange={1-6, 13-27}
% #2 = filenaam (moet in het directory ./progs staan)
\newcommand{\cpplstinputfragment}[2][]{%
	\lstinputfragment[%
	style=cppstyle,%  
	#1%
	]{#2}%
}

% De code in de environments cpplstfragment wordt in een codefragment geplaatst.
% #1 = (optional) extra key=value paren voor lstset
\lstnewenvironment{cpplstfragment}[1][]{%
	\lstset{
		style=cppstyle,%  
		aboveskip=.6\baselineskip,%
		#1%
	}%
}{%
}

% gebruik voor opmaak van warnings en error meldingen van de C++ compiler
\newcommand{\error}{\texttt}

% gebruikt voor in te vullen parameter in code
\newcommand{\param}[1]{\text{\color{cyan}{\textbf{#1}}}}

\begin{document}
\begin{cpplstShortInline}
% \begin{comment}
\tableofcontents
\section{Notes from A Tour if C++, Second Edition}
Aantekeningen die ik gemaakt heb bij het lezen van \emph{A Tour if C++, Second Edition} van Bjarne Stroustrup. 
De meeste aantekeningen zijn inmiddels verwerkt in het dictaat.
Wat overblijft vind je hier.

\subsection{The Basics}
\subsection{User-Defined Types}
\subsection{Modularity}
\begin{itemize}
	\item 
	Er was een voorstel voor \Cpp{}20 voor een mechanisme om contracten vast te leggen (invarianten, pre- en postcondities).
	Zie \url{http://www.open-std.org/JTC1/SC22/WG21/docs/papers/2015/n4415.pdf}. Maar dit voorstel heeft het niet gehaald.
	Wel in \Cpp{}26, zie \url{https://en.cppreference.com/cpp/language/contracts} en \url{https://en.cppreference.com/cpp/language/function#Function_contract_specifiers}.
\end{itemize}
\subsection{Classes}
\subsection{Essential Operations}
\subsection{Templates}
\begin{itemize}
	\item 
	Een lambda kan handig gebruikt worden om een groep statements te veranderen in een expressie.
	%TODO: kan dit gebruikt worden om een Breuk in te lezen in operator<<?
	Bijvoorbeeld bij een ingewikkelde initialisatie:
\begin{cpplstfragment}
	vector<int> v = [&] {
		switch (m) {
			case 0:
			return {0, 2, 4, 6, 8};
			case 1:
			return {1, 3, 5, 7, 9};
			default:
			return {};
		}
	};
\end{cpplstfragment}
	Is daar een toepassing voor in het dictaat?
\end{itemize}

\subsection{Strings and Regular Expressions}
\begin{itemize}
	\item 
	Bij het gebruik van regular expressions is het handig om een raw string literal te gebruiken |R("Hier heeft \ en " geen speciale betekenis")|.
	\item 
	Tip: gebruik |string| als returntype.
	Vertrouw op move semantics.
\end{itemize}
\subsection{Input and Output}
\subsection{Containers}
\begin{itemize}
	\item 
	|Save_vector| is een mooi voorbeeld van overerving.
	Zie \proglink{Save_vector.cpp}.
	Werkt niet polymorf omdat operator[] niet |virtual| gedefinieerd is.
	\item 
	Range checking kan soms ook aangezet worden als compiler optie.
	In gcc is dat de optie |-D_GLIBCXX_DEBUG|.
	\item 
	|count| bij |list| is $O(1)$ maar |count| bij |forward_list| is $O(n)$.	 
\end{itemize}
\subsection{Algorithms}
\begin{itemize}
	\item 
	Sinds \Cpp{}-17 bestaan er ook parallelle versies van de standaard algoritmen.
	%TODO kan ik niet compileren.
	Zie \proglink{sort_par.cpp}.
\end{itemize}
\subsection{Utilities}
\begin{itemize}
	\item 
	De standaard functie |forward| lijkt op de eenvoudigere |move| en kan gebruikt worden om een verzameling argumenten exact door te geven.
	Deze functie wordt onder andere gebruikt om |make_unique| te implementeren:
\begin{cpplstfragment}
template<typename T, typename ... Args>
unique_ptr<T> make_unique(Args&&... args) {
	return unique_ptr<T>{new T(std::forward<Args>(args)...)}
}
\end{cpplstfragment}
	\item 
	Sinds Cpp{}20 bevat de standaard library de class |span|.
	Dit is de combinatie van een pointer en een offset.
	\item 
	Andere utility classes: |bitset|, |pair|, |tuple|, |variant|, |optional| en |any|.
	|pair| en |tuple| zijn heterogeen alle andere containers zijn homogeen.
	|bitset| bevat een vast aantal bits.
	|variant| is een beter alternatief voor |union|.
	|optional| is een object dat mogelijk leeg is.
	|any| is een object dat elk willekeurig type kan bevatten.
	\item 
	Type deduction heeft de functie |make_pair| sinds \Cpp{}17 min of meer overbodig gemaakt.
	Zie \proglink{make_pair.cpp}.
	Hetzelfde geldt voor |make_tuple|.  
	\item 
	De \Cpp{}17 standaard library bevat een pool allocator.
	Deze allocator is sneller dan de standaard allocator, als je veel objecten van dezelfde grootte nodig hebt.
	Alle standaard containers hebben een extra template argument waarmee je de allocator kunt opgeven.
	\item 
	De include file |<chrono>| bevat utilities om met tijd om te gaan.
	Er zijn user defined suffixes gedefinieerd die je kunt gebruiken om je code leesbaarder te maken.
	Bijvoorbeeld:
\begin{cpplstfragment}
	this_thread::sleep(10ms+500us);
\end{cpplstfragment}
\end{itemize}

\subsection{Numerics}
\begin{itemize}
	\item 
	Als er een fout optreedt in een functie uit |<cmath>|, dan wordt dit gerapporteerd door de globale variabele |errno| te vullen met een foutcode.
\begin{cpplstfragment}
	errno = 0;
	double d {sqrt(-1)};
	if (errno == EDOM) {
		cerr << "Domain error for sqrt(-1)\n";
	}
\end{cpplstfragment}
	\item 
	De include file |<numeric>| definieert een aantal rekenkundige functies.
	Onder andere |inner_product|, |gcd| en |lcm|.
	|gcd| en |lcm| zijn toegevoegd in \Cpp{}17.
	\item 
	De minimale en maximale waardes van de ingebouwde types kunnen opgevraagd worden met functies die gedeclareerd zijn in |<limits>|.
	Bijvoorbeeld de maximale waarde voor een |int| is |numeric_limits<int>::max()|.
\end{itemize}
\subsection{Concurrency}
Kan eventueel nog aan dictaat worden toegevoegd.
Zie ook oude informatie uit RTSYST, zie \url{bd.eduweb.hhs.nl/rtsyst/studiew.htm}.
\subsection{History and Compatibility}
\begin{itemize}
	\item 
	De standaard library bevat ook de |ratio| voor \emph{compile-time rational arithmetic}.
	De teller en de noemer zijn template parameters.
	Een duidelijk use case ontbreekt.
\end{itemize}

\section{Notes from A Tour if C++, Third Edition}
Aantekeningen die ik gemaakt heb bij het lezen van \emph{A Tour if C++, Third Edition} van Bjarne Stroustrup. 
% De meeste aantekeningen zijn inmiddels verwerkt in het dictaat.
% Wat overblijft vind je hier.

\subsection{The Basics}
\begin{itemize}
	\item
	Een \emph{declaratie} is een instructie waarmee een entiteit in het programma wordt geïntroduceerd en waarvan het type wordt gespecificeerd:
	\begin{itemize}
		\item Een \emph{type} definieert een reeks mogelijke waarden en hoe die geïnterpreteerd moeten worden en een reeks bewerkingen (voor een object).
		\item Een \emph{object} is een stukje geheugen dat een waarde van een bepaald type bevat.
		\item Een \emph{waarde} is een reeks bits die volgens een bepaald type wordt geïnterpreteerd.
		\item Een \emph{variabele} is een object met een naam.
	\end{itemize}
	\item
	Gebruik |unsigned| uitsluitend voor bitmanipulatie; \href{https://isocpp.github.io/CppCoreGuidelines/CppCoreGuidelines#es101-use-unsigned-types-for-bit-manipulation}{CG: ES.101}.
\end{itemize}
\subsection{Modularity}
\subsection{Classes}
\subsection{Essential Operations}
\begin{itemize}
	\item
	Als een klasse X een destructor heeft die een niet-triviale taak uitvoert, zoals het vrijgeven van geheugen of een vergrendeling, heeft de klasse waarschijnlijk de volgende speciale memberfuncties nodig:
\begin{cpplstfragment}
class X {
public:
	X(Sometype); // constructor
	X(); // default constructor
	X(const X&); // copy constructor
	X(X&&); // move constructor
	X& operator=(const X&); // copy assignment
	X& operator=(X&&); // move assignment
	~X(); // destructor
	// ...
};
\end{cpplstfragment}
	\item
	Behalve de \squote{gewone} constructor worden deze speciale memberfuncties door de compiler gegenereerd wanneer dat nodig is.
	Als je expliciet wilt aangeven dat er standaardimplementaties moeten worden gegenereerd kun je de syntax |=default| gebruiken.
	Als je voor bepaalde speciale memberfuncties expliciet bent, worden andere speciale memberfuncties niet gegenereerd.
	Als je expliciet wilt aangeven dat er geen standaardimplementatie gegenereerd mag worden, kun je de syntax |=delete| gebruiken.
\end{itemize}
\subsection{Templates}
\begin{itemize}
	\item
	Het zelfgedefinieerde type |Array| kan als \dquote{\USE{constrained template}} (to constrain = beperken) worden gedefinieerd, zodat het alleen gebruikt kan worden voor types die voldoen aan bepaalde eisen.
	Een constrained template heeft een zogenoemd \dquote{\USE{\emph{concept}}} als template parameter.
\begin{cpplstfragment}
template<copyable T> class Array {
	// ...
};
\end{cpplstfragment}
\end{itemize}
\subsection{Concepts and Generic Programming}
\begin{itemize}
	\item
	De concepten die voor een template zijn gespecificeerd, worden gebruikt om argumenten te controleren op het moment dat het template wordt gebruikt.
	Ze worden niet gebruikt om het gebruik van de parameters in de definitie van het template te controleren.
	\item
	\dquote{\USE{Fold expressions}} zijn een handige manier om een parameter pack te verwerken.
	Zie \proglink{fold_expression.cpp}.
	\item
	|std::forward<>| kan gebruikt worden om een parameter pack exact door te geven aan een andere functie, bijvoorbeeld in de implementatie van |make_unique|:
\begin{cpplstfragment}
template<typename T, typename ... Args>
unique_ptr<T> make_unique(Args&&... args) {
	return unique_ptr<T>{new T(std::forward<Args>(args)...)};
}
\end{cpplstfragment}
\end{itemize}

\section{Notes from Effective Modern C++}
Aantekeningen die ik gemaakt heb bij het lezen van \emph{Effective Modern C++} van Scott Meyers. 

\section{Notes from Professional C++, 6th Edition}
Aantekeningen die ik gemaakt heb bij het lezen van \emph{Professional C++, 6th Edition} van Gregoire.
Programmacode: \url{https://github.com/Professional-CPP/edition-6}.
\subsection{C++ Crash Course}
\begin{itemize}
	\item
	|endl| flusht het output buffer |"\n"| niet.
	\item
	In een header file moet je nooit een using directive gebruiken, omdat dit kan leiden tot naamconflicten in de code die deze header file include.
	In een module interface file kun je wel een using directive gebruiken, zolang je die maar niet exporteert.
	\item
	Types |char8_t|, |char16_t| en |char32_t| zijn bedoeld voor het opslaan van respectievelijk UTF-8, UTF-16 en UTF-32 karakters.
	NOTE: niet opgenomen omdat een duidelijk voorbeeld ontbreekt.
	\item
	Attributes.
	Eventueel nog uitbreiden met:
	\begin{itemize}
		\item |[[deprecated]]|.
		\item |[[likely]]| and |[[unlikely]]|.
		\item |[[assume]]|.
	\end{itemize}
	\item
	De term \dquote{functiesignatuur} (Engels: function signature) wordt gebruikt om de combinatie van de functienaam en de parameterlijst aan te duiden, maar zonder het return type.
	\item
	Je kunt de naam van een functie opvragen via |__func__|. 
	Deze naam is een string literal die de naam van de functie bevat waarin deze wordt gebruikt.
	Zie \proglink{__func__.cpp}
	\item
	Met |std::size()| kun je het aantal elementen in een C-array opvragen.
	Zie \proglink{size.cpp}.
	\item
	Met class template argument deduction (CTAD) wordt de template argument van een class template automatisch afgeleid uit de constructor argumenten.
	Dat kan uiteraard alleen als de class template m.b.v. een constructor met een argument dat door de compiler kan worden bepaald wordt aangeroepen.
	Zie \proglink{CTAD.cpp}.
	\item
	Hoewel je een functie zou kunnen schrijven die een |pair| retourneert, wordt aangeraden om een kleine |struct| te maken die de twee waarden bevat en die vanuit de functie te retourneren.
	\item
	|optional<T>| kan een waarde of niets bevatten. 
	Je kunt controleren of een |optional| een waarde bevat met behulp van de |has_value()| functie of door de |optional| te gebruiken in een boolean context.
	De waarde van een |optional| kan worden opgevraagd met behulp van de |value()| functie of door de |optional| te dereferencen met behulp van de |*| operator.
	Zie \proglink{optional.cpp}.
	\item
	M.b.v. aangewezen initialisatie (Engels \USE{designated initializers}) kun je datamembers initialiseren m.b.v. hun naam.
	Deze syntax is vooral handig bij het initialiseren van structuren met veel datamembers, omdat het de leesbaarheid van de code verbetert en het risico op fouten bij het initialiseren van de verkeerde datamember vermindert.
	Zie \proglink{designated_initializers.cpp}.
	De volgorde van de designated initializers moet overeenkomen met de volgorde van de datamembers in de struct.
	Een designated initializer kan overgeslagen worden als het datamember een defaultwaarde heeft.
	\item
	De heap wordt ook wel de free store genoemd.
	\item
	|as_const(obj)| geeft een const reference terug naar |obj|.
	\item
	Een type alias is beter dan een |typedef| omdat ze beter leesbaar zijn en omdat ze kunnen worden gebruikt met templates.
	Voorbeeld:
\begin{cpplstfragment}
	template<typename T>
	using Vec = vector<T>;
	Vec<int> v; // is hetzelfde als vector<int> v;
	using map_t = map<string, int>;
	map_t m; // is hetzelfde als map<string, int> m;
\end{cpplstfragment}
	\item
	|auto| stript qualifiers (|const| en |volatile|) en references van het type dat wordt afgeleid weg.
	Dit kan tot onverwachte kopietjes leiden.
	|decltype| doet dat niet, maar behoudt alle qualifiers en references.
\end{itemize}
\subsection{Working with Strings and String Views}
\begin{itemize}
	\item
	\USE{\emph{raw string literals}}, zie \proglink{raw_string_literal.cpp}.
	De delimiter kan tussen de 1 en 16 karakters lang zijn.
	\item
	Alles wat in een \USE{\emph{inline namespace}} wordt gedefinieerd is automatisch ook beschikbaar in de bovenliggende namespace, zie \proglink{inline_namespace.cpp}.
	\item
	CTAD met |vector| met strings.
	Je moet oppassen wanneer je CTAD gebruikt voor een |vector| met strings.
	Een string literal is namelijk niet van het type |string| maar van het type |const char*| of |const char[]|.
	Zie \proglink{CTAD_vector_string.cpp}.
	\item
	De volgende functies zijn beschikbaar om numerieke waarden om te zetten in een |string|, waarbij het type \code[mathescape=true]{$\param{T}$} (|unsigned|) |int|, (|unsigned|) |long|, (|unsigned|) |long long|, |float|, |double| of |long double| kan zijn.
\begin{cpplstfragment}[mathescape=true]
string to_string($\param{T}$ value);
\end{cpplstfragment}
	Zie \proglink{to_string.cpp}.
	\item
	De volgende functies zijn beschikbaar om een |string| om te zetten in een numerieke waarde.
\begin{cpplstfragment}[mathescape=true]
$\param{T}$ sto$\param{X}$(const string& str, size_t *pos = nullptr, int base = 10);
\end{cpplstfragment}
	Waarbij \code[mathescape=true]{$\param{X}$} |i|, |ui|, |l|, |ul|, |ll|, |ull|, |f|, |d| of |ld| kan zijn.
	\code[mathescape=true]{$\param{T}$} is dan respectievelijk |int|, |unsigned int|, |long|, |unsigned long|, |long long|, |unsigned long long|, |float|, |double| of |long double|.
	De functies die een floating point waarde teruggeven, hebben geen parameter |base|.
	|pos| is een output parameter. Na afloop bevat |*pos| de index van het eerste niet-omgezette teken. 
	|base| geeft het grondtal aan dat tijdens de omzetting moet worden gebruikt.
	Als |base| 0 is, wordt het grondtal automatisch bepaald op basis van de inhoud van de string:
	\begin{itemize}
		\item als de string begint met |0|, wordt het grondtal 8;
		\item als de string begint met |0x|, wordt het grondtal 16;
		\item anders wordt het grondtal 10\footnote{%
			Je zou misschien verwachten dat als de string begint met \code{0b}, het grondtal 2 zou worden, maar dat is (helaas) niet het geval in \Cpp{}23.
			Zie \url{https://en.cppreference.com/w/cpp/string/basic_string/stol}.
		}.
	\end{itemize}
	Deze functies slaan voorafgaande spaties over.
	Er wordt een genereren een |invalid_argument| exception gegooid als er geen conversie kan worden uitgevoerd.
	Er wordt een |out_of_range| exception gegooid als de geconverteerde waarde buiten het bereik van het return type valt.
	Zie \proglink{sto_numeriek.cpp}.
	\item
	Er zijn low level conversie functies beschikbaar (|from_chars|\footnote{%
		Zie \url{https://en.cppreference.com/cpp/utility/from_chars}.
	} en |to_chars|\footnote{%
		Zie \url{https://en.cppreference.com/cpp/utility/to_chars}.
	}) om numerieke gegevens te serialiseren naar en te deserialiseren vanuit formaten zoals JSON, XML, enzovoort..
	\item
	Je kunt een string literal afsluiten met de |sv| literal om deze als |string_view| te interpreteren:
\begin{cpplstfragment}
	auto hallo = "Hallo Wereld!"sv; // hallo is van het type string_view
\end{cpplstfragment}
	\item
	In de format string van een |format|, |print| of |println| functieaanroep wordt een |{{| geïnterpreteerd als |{| en een |}}| als |}|, zodat je deze tekens in de output kunt opnemen zonder dat ze worden geïnterpreteerd als format specifiers.
	\item
	Bla
	\item
	Bla
	\item
	Bla
	\item
	Bla
	\item
	Bla
	\item
	Bla
	\item
	Bla
\end{itemize}

\section{Testen commando om een breuk te printen}
Dit is een regel zonder bereuken erin, dat moet natuurlijk ook gewoon kunnen en kan gebruikt worden om te testen of de regelafstand niet teveel veranderd.
De waarde van een half is \myfrac{1}{2}.
De waarde van een zeven vijftiende is \myfrac{7}{5}.
De waarde van eenendertig honderdste is \myfrac{31}{100}.
Dit is een regel zonder bereuken erin, dat moet natuurlijk ook gewoon kunnen en kan gebruikt worden om te testen of de regelafstand niet teveel veranderd.
De waarde van een half is \myfrac{1}{2}.
De waarde van een zeven vijftiende is \myfrac{7}{5}.
De waarde van eenendertig honderdste is \myfrac{31}{100}.
Dit is een regel zonder bereuken erin, dat moet natuurlijk ook gewoon kunnen en kan gebruikt worden om te testen of de regelafstand niet teveel veranderd.
\section{Testen van keywords highlighting}
%test keywords highlighting
\begin{cpplstfragment}
alignas (C++11)
alignof (C++11)
and
and_eq
asm
atomic_cancel (TM TS)
atomic_commit (TM TS)
atomic_noexcept (TM TS)
auto (1) (3) (4) (5)
bitand
bitor
bool
break
case
catch
char
char8_t (C++20)
char16_t (C++11)
char32_t (C++11)
class (1)
compl
concept (C++20)
const
consteval (C++20) (5)
constexpr (C++11) (3)
constinit (C++20)
const_cast
continue
contract_assert (C++26)
co_await (C++20)
co_return (C++20)
co_yield (C++20)
decltype (C++11) (2)
default (1)
delete (1)
do
double
dynamic_cast
else
enum (1)
explicit
export (1) (4)
extern (1)
false
float
for (1)
friend
goto
if (3) (5)
inline (1) (3)
int (1)
long
mutable (1)
namespace
new
noexcept (C++11)
not
not_eq
nullptr (C++11)
operator (1)
or
or_eq
private (4)
protected
public
reflexpr (reflection TS)
register (3)
reinterpret_cast
requires (C++20)
return
short
signed
sizeof (1)
static
static_assert (C++11)
static_cast
struct (1)
switch
synchronized (TM TS)
template
this (5)
thread_local (C++11)
throw (3) (4)
true
try
typedef
typeid
typename (3) (4)
union
unsigned
using (1) (4)
virtual
void
volatile
wchar_t
while
xor
xor_eq 
\end{cpplstfragment}

(1) — meaning changed or new meaning added in C++11.
(2) — new meaning added in C++14.
(3) — meaning changed or new meaning added in C++17.
(4) — meaning changed or new meaning added in C++20.
(5) — new meaning added in C++23. 

In addition to keywords, there are identifiers with special meaning, which may be used as names of objects or functions, but have special meaning in certain contexts:

\begin{cpplstfragment}
final (C++11)
override (C++11)
transaction_safe (TM TS)
transaction_safe_dynamic (TM TS)
import (C++20)
module (C++20)
pre (C++26)
post (C++26)
trivially_relocatable_if_eligible (C++26)
replaceable_if_eligible (C++26)
\end{cpplstfragment}
% \end{comment}

\section{Testcode}


% \begin{comment}
\section{Bruikbare tekst}
Het \USE{keyword} |typename| is nodig om de compiler duidelijk te maken dat \hcode{iterator_\\-traits<I>::\\-iterator_\\-category} een typenaam is.
Zie eventueel \url{http://en.cppreference.com/w/cpp/keyword/typename}.


% \end{comment}
\end{cpplstShortInline}
\end{document}